% Rapport de sécurité — Tropic Line (ShopCart), projet DevSecOps ING-5-SSIR, TEK-UP.
% Compiler : pdflatex rapport-securite.tex (deux fois, pour la table des matières et les liens)
% ou importer ce fichier dans Overleaf (compilateur pdfLaTeX).
\documentclass[11pt,a4paper]{article}

\usepackage[T1]{fontenc}
\usepackage[utf8]{inputenc}
\usepackage{lmodern}
\usepackage[french]{babel}
\usepackage[a4paper,top=2cm,bottom=2.2cm,left=2.1cm,right=2.1cm]{geometry}
\usepackage{microtype}
\usepackage{xcolor}
\usepackage{booktabs,tabularx,array,longtable}
\usepackage{enumitem}
\usepackage{fancyhdr}
\usepackage{tikz}
\usetikzlibrary{positioning,arrows.meta,fit,babel}
\usepackage[hidelinks]{hyperref}
\usepackage{caption}
\usepackage{float}

% ---------------------------------------------------------------- à compléter par l'équipe
\newcommand{\Equipe}{Nom Prénom 1, Nom Prénom 2, Nom Prénom 3, Nom Prénom 4}
\newcommand{\Depot}{\url{https://github.com/zied458/ShopCart}}
% ----------------------------------------------------------------

\definecolor{ink}{HTML}{21170F}
\definecolor{cherry}{HTML}{A3172C}
\definecolor{ok}{HTML}{1E7B34}
\definecolor{ko}{HTML}{B3261E}
\definecolor{soft}{HTML}{F3F1EE}
\newcommand{\ok}{\textcolor{ok}{\textbf{OK}}}
\newcommand{\ko}{\textcolor{ko}{\textbf{échec}}}
\newcommand{\code}[1]{\texttt{\small #1}}

\setlist{nosep,leftmargin=*}
\setlength{\parskip}{0.35em}
\setlength{\parindent}{0pt}
\renewcommand{\arraystretch}{1.18}
\captionsetup{font=small,labelfont=bf,skip=4pt}
\newcolumntype{Y}{>{\raggedright\arraybackslash}X}
\newcolumntype{L}[1]{>{\raggedright\arraybackslash}p{#1}}

\pagestyle{fancy}
\fancyhf{}
\lhead{\small Tropic Line --- Rapport de sécurité}
\rhead{\small DevSecOps, ING-5-SSIR}
\cfoot{\small \thepage}
\renewcommand{\headrulewidth}{0.3pt}

\hypersetup{pdftitle={Rapport de sécurité --- Tropic Line (ShopCart)},pdfauthor={Équipe Tropic Line}}

\begin{document}

% ================================================================ titre
\begin{center}
  {\color{cherry}\rule{\linewidth}{1.2pt}}\\[0.6em]
  {\LARGE\bfseries Rapport de sécurité --- Tropic Line (ShopCart)}\\[0.4em]
  {\large Sécuriser une boutique en ligne sur tout son cycle de vie}\\[0.7em]
  {\small TEK-UP University --- ING-5-SSIR --- Module DevSecOps --- Enseignant : Oussema Mansouri}\\[0.3em]
  {\small Équipe : \Equipe}\\[0.2em]
  {\small Dépôt : \Depot\quad --- \quad Novembre 2026}\\[0.4em]
  {\color{cherry}\rule{\linewidth}{1.2pt}}
\end{center}

\noindent\colorbox{soft}{\parbox{\dimexpr\linewidth-2\fboxsep}{\small
\textbf{Résumé.} Tropic Line est une boutique de cafés de spécialité (React, Node.js/Express, PostgreSQL)
construite avec \textbf{huit failles volontaires}, documentées avant toute correction. Une chaîne DevSecOps
a ensuite été mise en place : images conteneur minimales et sans root, pipeline GitHub Actions à étape de
sécurité bloquante (Gitleaks, Semgrep, Trivy), secrets fournis à l'exécution et révoqués après la fuite,
déploiement Kubernetes durci (Pod Security \og restricted \fg{}, NetworkPolicies en refus par défaut,
Kyverno), scan dynamique authentifié OWASP ZAP, détection à l'exécution et plan de réponse à incident.
Les huit failles sont corrigées, chacune dans son propre commit, et protégées par un test de régression :
\textbf{13 tests sur 18 échouaient} sur la version vulnérable, \textbf{19 sur 19 réussissent} aujourd'hui.}}

\vspace{0.2em}
{\small\tableofcontents}

% ================================================================
\section{Contexte et périmètre}

\textbf{L'application.} Une boutique en ligne pour un torréfacteur fictif : six cafés d'origine (Inde,
Brésil, Éthiopie\ldots) et du matériel d'extraction. Les clients créent un compte, commandent, suivent leurs
commandes et laissent des avis ; un administrateur gère commandes, prix, stocks et clients.

\textbf{Données sensibles.} Comptes (e-mail, empreinte bcrypt du mot de passe, rôle), coordonnées de
livraison (nom, téléphone, adresse), historique de commandes, et la clé de signature des jetons de session.

\textbf{Technologies.} Boutique : React 19, TypeScript, Vite, Tailwind CSS, servie par nginx. API : Node.js 24,
Express 5, jetons JWT (HS256), bcrypt. Base : PostgreSQL 17 (tables \code{users}, \code{products},
\code{orders}, \code{order\_items}, \code{reviews}, \code{security\_events}). Livraison : Docker,
GitHub Actions, Kubernetes (kind), Kyverno, OWASP ZAP.

\begin{figure}[H]
\centering
\begin{tikzpicture}[
  node distance=0.6cm and 1.15cm,
  box/.style={draw=ink,rounded corners=2pt,minimum height=0.85cm,align=center,font=\small,fill=white},
  arr/.style={-{Stealth[length=2mm]},thick,color=ink},
  lbl/.style={font=\scriptsize,color=ink!70}]
  \node[box] (user) {Navigateur\\client / admin};
  \node[box,right=of user] (web) {\textbf{web}\\nginx (UID 101)\\×2};
  \node[box,right=of web] (api) {\textbf{api}\\Node.js (UID 65532)\\×2};
  \node[box,right=of api] (db) {\textbf{db}\\PostgreSQL (UID 70)};
  \node[box,below=of api,font=\scriptsize] (sec) {Secret \code{tropicline-secrets}\\(fichiers en lecture seule)};
  \draw[arr] (user) -- node[above=1pt,lbl]{8090} (web);
  \draw[arr] (web) -- node[above=1pt,lbl]{3000} (api);
  \draw[arr] (api) -- node[above=1pt,lbl]{5432} (db);
  \draw[arr,dashed] (sec) -- (api);
  \node[draw=cherry,dashed,rounded corners,inner sep=7pt,fit=(web)(api)(db)(sec),
        label={[font=\scriptsize,color=cherry]below:namespace \code{tropicline} --- Pod Security restricted, refus réseau par défaut, Kyverno}] {};
\end{tikzpicture}
\caption{Architecture déployée et frontières de confiance. Seuls les flux fléchés sont autorisés.}
\end{figure}

\textbf{Éthique.} Toutes les attaques ont été menées uniquement contre nos propres instances locales
(Docker Compose et cluster kind), avec des comptes et des secrets de test.

% ================================================================
\section{Modèle de menace (STRIDE)}

\textbf{Actifs à protéger} : données personnelles des clients ; comptes, en particulier administrateur ;
clé de signature des jetons ; intégrité des prix, des stocks et des commandes ; disponibilité de la
boutique ; chaîne de livraison (dépôt, pipeline, images, cluster).

\textbf{Frontières de confiance} : navigateur $\rightarrow$ nginx (Internet, non fiable) ; nginx
$\rightarrow$ API ; API $\rightarrow$ base ; dépôt $\rightarrow$ CI $\rightarrow$ registre
$\rightarrow$ cluster. Tout ce qui vient du navigateur (corps, paramètres, en-têtes, jeton) est contrôlé
par l'attaquant.

\begin{table}[H]
\small
\begin{tabularx}{\linewidth}{@{}L{2.1cm}YL{1.1cm}Y@{}}
\toprule
\textbf{STRIDE} & \textbf{Menace sur Tropic Line} & \textbf{Faille} & \textbf{Contrôle en place} \\
\midrule
\textbf{S}poofing & Fabriquer un jeton admin avec la clé lue dans le dépôt ; deviner un mot de passe & V6 & secret à l'exécution, rotation ; verrouillage après 5 échecs \\
\textbf{T}ampering & Modifier le prix payé ; polluer \code{Object.prototype} ; image ou action CI altérée & V2, V7 & prix lu en base ; champs nommés ; images et actions figées par empreinte, registres autorisés (Kyverno) \\
\textbf{R}epudiation & Agir sans laisser de trace (aucun journal à l'origine) & A09 & journaux JSON corrélés, audit des actions admin \\
\textbf{I}nformation disclosure & Injection SQL ; lire la commande d'autrui ; pile d'appels et erreurs SQL renvoyées & V1, V3, V8 & requêtes paramétrées ; contrôle de propriété ; erreurs génériques \\
\textbf{D}enial of service & Épuiser le serveur ; verrouiller le compte d'autrui & --- & limites CPU/mémoire, quotas, corps $\leq$ 100 Ko ; verrouillage par couple adresse + compte \\
\textbf{E}levation of privilege & S'inscrire en admin ; script volant la session d'un admin ; sortie du conteneur & V4, V5 & rôle fixé et vérifié en base ; rendu texte + CSP ; non-root, sans capacités, lecture seule \\
\bottomrule
\end{tabularx}
\caption{Analyse STRIDE : menaces, failles du projet correspondantes et contrôles.}
\end{table}

% ================================================================
\section{Failles initiales}

La première version (commit \code{c70fc4a}) contient huit failles volontaires, chacune marquée \code{Vn}
dans le code et documentée \emph{avant} correction (emplacement, cause, impact) dans
\code{docs/vulnerabilities.md}. Chaque impact a été démontré sur notre instance.

\begin{table}[H]
\small
\begin{tabularx}{\linewidth}{@{}L{0.6cm}L{3.2cm}L{3.7cm}YL{1.3cm}@{}}
\toprule
\textbf{\#} & \textbf{Faille} & \textbf{Emplacement} & \textbf{Impact démontré} & \textbf{Gravité} \\
\midrule
V1 & Injection SQL & \code{/api/products?search=} & une requête \code{UNION SELECT} renvoie 7 comptes avec e-mail et empreinte du mot de passe & critique \\
V2 & Prix fourni par le client & \code{POST /api/orders} & moulin à 165 TND acheté 0,001 TND & élevée \\
V3 & IDOR sur les commandes & \code{GET /api/orders/:id} & un client lit nom, téléphone, adresse d'un autre & élevée \\
V4 & Rôle choisi à l'inscription & \code{/api/auth/register} & \code{"role":"admin"} crée un administrateur & critique \\
V5 & XSS stocké & avis produit (API + page) & un script dans un avis s'exécute chez chaque visiteur & élevée \\
V6 & Secret dans le code & \code{config.js}, Compose & jeton admin fabriqué : 6 fiches clients lues (HTTP 200) & critique \\
V7 & Dépendance vulnérable & lodash 4.17.11 & pollution de prototype (CVE-2019-10744) modifiant le serveur & critique \\
V8 & Mauvaise configuration & \code{app.js} & CORS \code{*}, erreurs SQL et piles renvoyées, \code{X-Powered-By} & moyenne \\
\bottomrule
\end{tabularx}
\caption{Failles de la version initiale (gravité qualitative : impact × facilité d'exploitation).}
\end{table}

% ================================================================
\section{Cartographie de référence}

Chaque faille est rattachée à l'OWASP Top 10 (édition 2025), à la CWE et à STRIDE, avec les contrôles qui
l'ont vue.

\begin{table}[H]
\small
\begin{tabularx}{\linewidth}{@{}L{0.55cm}L{3.1cm}L{2.2cm}L{1.35cm}YYL{0.8cm}@{}}
\toprule
\textbf{\#} & \textbf{OWASP Top 10:2025} & \textbf{CWE} & \textbf{STRIDE} & \textbf{Analyse statique (CI)} & \textbf{ZAP (exécution)} & \textbf{Test} \\
\midrule
V1 & A05 Injection & CWE-89 & I & Semgrep (règle d'équipe) & \textbf{High} confirmé & 2 \\
V2 & A06 Insecure Design & CWE-602 & T & --- & --- & 1 \\
V3 & A01 Broken Access Control & CWE-639 & I & --- & --- & 2 \\
V4 & A01 Broken Access Control & CWE-915, 269 & E & Semgrep (règle d'équipe) & --- & 1 \\
V5 & A05 Injection & CWE-79 & E & Semgrep (règle d'équipe) & confirmé (persistant) & 3 \\
V6 & A04 Cryptographic Failures & CWE-321, 798 & S & Gitleaks, Semgrep & --- (corrigée avant) & --- \\
V7 & A03 Supply Chain Failures & CWE-1395, 1321 & T & Trivy (dépendances, image) & --- & 2 \\
V8 & A02 Misconfiguration, A10 & CWE-209, 942, 693 & I & Semgrep (2 avertissements) & confirmé (CORS, en-têtes) & 4 \\
\bottomrule
\end{tabularx}
\caption{Cartographie des failles. \og --- \fg{} : non détectée par ce moyen.}
\end{table}

\begin{table}[H]
\small
\begin{tabularx}{\linewidth}{@{}L{4.2cm}YL{4.6cm}@{}}
\toprule
\textbf{Contrôle} & \textbf{Ce qu'il couvre} & \textbf{Référence} \\
\midrule
Gitleaks (tout l'historique) & secrets dans le code et l'historique & A04, CWE-798 \\
Semgrep (\code{p/default} + règles d'équipe) & injection, XSS, rôle lu dans la requête, secrets & A05, A01, A02 \\
Trivy \code{fs} et \code{image} & CVE des dépendances et des paquets système, secrets dans l'image & A03, CWE-1395 \\
Trivy \code{config}, Pod Security, Kyverno & privilèges, root, capacités, registres, limites & A02, \og restricted \fg{} (PSS) \\
NetworkPolicies, comptes de service & mouvement latéral, exfiltration, accès à l'API Kubernetes & moindre privilège \\
ZAP authentifié & injection et configuration à l'exécution & A05, A02 \\
Journaux, règles de détection & attaques en cours, traçabilité & A09 \\
\bottomrule
\end{tabularx}
\caption{Contrôles de la chaîne DevSecOps et références couvertes.}
\end{table}

% ================================================================
\section{Correctifs}

Méthode : (1) écrire un test de régression du point de vue de l'attaquant, exécuté contre l'application
déployée ; (2) vérifier qu'il échoue sur la version vulnérable ; (3) corriger la \textbf{cause} ; (4)
vérifier test, Semgrep, Trivy et ZAP ; (5) un commit par faille. Les tests tournent ensuite dans la CI
après chaque déploiement : une faille réintroduite casse le pipeline.

\begin{table}[H]
\small
\begin{tabularx}{\linewidth}{@{}L{0.6cm}YL{4.3cm}@{}}
\toprule
\textbf{\#} & \textbf{Correctif} & \textbf{Preuve} \\
\midrule
V1 & Requête paramétrée (\code{ILIKE \$1}) : le texte est une donnée, jamais du code ; \code{\%} et \code{\_} échappés ; 100 caractères au plus ; tri choisi dans une liste fixe. & tests 2/2 ; Semgrep ; ZAP : High disparu \\
V2 & Prix lu dans \code{products}, prix du client ignoré (et plus envoyé) ; ligne verrouillée (\code{FOR UPDATE}) contre la double vente du dernier sachet. & test : commande à 0,001 facturée 18 TND \\
V3 & La commande doit appartenir à l'appelant ; sinon \textbf{404}, comme une commande inexistante (ne confirme pas les numéros utilisés). & tests 2/2 \\
V4 & Rôle toujours \code{customer} ; le rôle admin est relu \textbf{en base} à chaque requête, pas dans le jeton. Remédiation : 3 comptes déjà élevés rétrogradés. & test : compte créé \code{customer}, route admin 403 \\
V5 & Avis affiché comme \textbf{texte} (échappé par React) ; défense en profondeur : CSP \code{script-src 'self'} sans script en ligne ; 1\,000 caractères au plus. & tests 3/3 ; navigateur : aucun script exécuté, aucune violation CSP \\
V6 & Secrets fournis en fichiers à l'exécution, refus de démarrer sans eux ; \textbf{révoqués d'abord} (nouvelle clé, nouveau mot de passe) ; purge de l'historique préparée. & jeton forgé : 200 $\rightarrow$ 401 \\
V7 & lodash \textbf{supprimé} ; la mise à jour du profil copie des champs nommés, type et longueur vérifiés. & test de pollution ; Trivy : 1 critique + 3 élevées $\rightarrow$ 0 \\
V8 & CORS supprimé (même origine via nginx), \code{X-Powered-By} désactivé, en-têtes de sécurité, erreurs génériques avec identifiant de corrélation, détail dans le journal seulement. & tests 4/4 ; ZAP : 9 alertes disparues \\
\bottomrule
\end{tabularx}
\caption{Correctifs (un commit par faille) et preuves.}
\end{table}

\textbf{Choix notables.} Pour V1, filtrer les apostrophes ou le mot \code{UNION} se contourne et casse des
recherches légitimes ; la requête paramétrée supprime la cause. Pour V3, répondre 404 plutôt que 403 évite
de révéler quels numéros existent. Pour V7, mettre lodash à jour aurait corrigé les CVE connues, mais
fusionner tout un corps de requête reste une affectation de masse ; supprimer la dépendance retire aussi
une surface d'attaque. Pour V4, corriger le code ne suffit pas : les comptes déjà élevés ont été
rétrogradés.

\begin{table}[H]
\small
\begin{tabularx}{\linewidth}{@{}YYY@{}}
\toprule
\textbf{Contrôle} & \textbf{Version vulnérable} & \textbf{Version corrigée} \\
\midrule
Tests de régression & 13 échecs sur 18 & \textbf{19 / 19} (lockout inclus) \\
Semgrep & 4 \code{ERROR} (V1, V4, V5, V6) & \textbf{0 ERROR} (2 avertissements, faux positifs) \\
Trivy, dépendances de l'API & 1 critique, 3 élevées & \textbf{0} \\
ZAP authentifié & 1 High, 5 Medium, 3 Low & \textbf{0 High réelle}, 11 alertes sur 13 disparues \\
Gitleaks & 3 secrets & \textbf{0} (révoqués, puis historique purgé) \\
\bottomrule
\end{tabularx}
\caption{Avant / après.}
\end{table}

% ================================================================
\section{Chaîne DevSecOps}

\subsection{Conteneurs}
Images multi-étapes ; API sur base \emph{distroless} (ni shell ni gestionnaire de paquets), utilisateur
65532 ; boutique sur nginx non privilégié, utilisateur 101, paquets mis à jour et \code{curl} retiré. Images
de base figées par empreinte, \code{.dockerignore} en liste blanche (pas de \code{.git}, pas de secret), code
appartenant à root donc non modifiable par le processus. Exécution en lecture seule, sans capacité Linux,
\code{no-new-privileges}.

\subsection{Pipeline et politique de blocage}
GitHub Actions, outils dans des images à version figée et actions figées par empreinte de commit. Chaque
barrière publie un tableau lisible sur la page de l'exécution, et un job final les résume.

\begin{table}[H]
\small
\begin{tabularx}{\linewidth}{@{}L{3.6cm}L{4.2cm}Y@{}}
\toprule
\textbf{Barrière} & \textbf{Bloque sur} & \textbf{Justification} \\
\midrule
Secrets (Gitleaks) & toute détection, tout l'historique & un secret poussé est compromis ; valeurs caviardées \\
Code (Semgrep) & \code{ERROR} & failles exploitables ; \code{WARNING} trop dépendant du contexte \\
Dépendances, images (Trivy) & \code{HIGH}/\code{CRITICAL} avec correctif ; image en root & corriger coûte une mise à jour ; sans correctif, bloquer figerait tout \\
Fichiers de déploiement & \code{HIGH}/\code{CRITICAL} & privilèges menant au nœud \\
Déploiement de test & 22 contrôles de durcissement & contrôles binaires \\
Tests, ZAP, détection & un test, une alerte \code{High}, une règle muette & preuve à l'exécution \\
\bottomrule
\end{tabularx}
\caption{Politique de blocage (détail et exceptions datées : \code{docs/pipeline-policy.md}).}
\end{table}

Le premier passage sur GitHub a échoué comme prévu (secrets, code, dépendances, image API) et n'a rien
publié. Une exception n'est acceptée qu'avec une justification, un responsable et une échéance ; un vrai
secret n'est jamais mis en exception.

\subsection{Secrets : révoquer d'abord, purger ensuite}
Les secrets sont des fichiers montés (secrets Docker, Secret Kubernetes en lecture seule, clé par clé),
générés aléatoirement au déploiement ; l'API refuse de démarrer sans eux. Après la fuite V6 : (1)
\textbf{révocation} --- nouveau mot de passe (\code{ALTER USER}), nouvelle clé de signature, redémarrage :
le jeton forgé passe de \textbf{200 à 401}, l'ancien mot de passe est refusé, les clients se reconnectent ;
(2) \textbf{purge} --- \code{git filter-repo} remplace les valeurs dans tous les commits d'un miroir, vérifié
par Gitleaks (\textbf{3 $\rightarrow$ 0}), puis envoi forcé : l'étape \og Secrets \fg{} du pipeline
repasse au vert. Purger sans révoquer laisserait le secret valide dans tous les
clones déjà faits.

\subsection{Déploiement Kubernetes durci}
Namespace en Pod Security \og restricted \fg{} ; un compte de service par composant, \textbf{aucun droit} sur
l'API Kubernetes et aucun jeton monté ; \code{securityContext} : non-root, lecture seule, sans capacité,
seccomp, limites ; NetworkPolicies en \textbf{refus par défaut} (web $\rightarrow$ api $\rightarrow$ db,
DNS seulement si nécessaire, aucune sortie vers Internet) ; quatre politiques Kyverno en mode \code{Deny}
(durcissement, limites, registres et versions explicites, secrets en fichiers). \code{verify.sh} prouve
\textbf{22 contrôles sur le cluster en marche} : par exemple, un pod root ou une image \code{nginx:latest}
est refusé, la boutique ne joint pas la base, l'API ne joint pas Internet.

\subsection{Test dynamique authentifié}
ZAP 2.17 se connecte comme client et envoie le jeton sur chaque requête (preuve : \code{/api/auth/me} doit
répondre 200), importe la description OpenAPI de l'API et explore la boutique. Avant correction, il
confirme \textbf{V1} (injection SQL, High), \textbf{V5} et \textbf{V8}. Après correction, il ne reste
aucune alerte High réelle ; un faux positif \og Path Traversal \fg{} (confiance faible) est déclaré pour
une seule URL et un seul paramètre, avec sa raison.

\subsection{Détection à l'exécution et réponse}
L'API écrit une ligne JSON par requête et par événement de sécurité, avec l'identifiant de corrélation de
nginx ; e-mails masqués, ni mot de passe ni jeton. Trois règles sur 5 minutes : \textbf{brute force}
(5 échecs sur un compte depuis une adresse $\rightarrow$ verrouillage 429), \textbf{sonde d'injection}
(3 motifs SQL ou script), \textbf{sonde d'accès} (10 refus ou commandes inconnues). Les compteurs sont
partagés en base entre les répliques ; des métriques Prometheus restent internes. La CI vérifie que les
attaques du job (tests, ZAP) déclenchent les trois règles. Le plan de réponse
(\code{docs/incident-response.md}) définit rôles, gravités P1--P3 et cinq scénarios avec commandes, dont
la mise en quarantaine d'un pod ; il a été exercé sur le cluster.

% ================================================================
\section{Réflexion critique}

\textbf{Ce que les outils voient, et ce qu'ils ne voient pas.}
\begin{itemize}
  \item Les règles publiques de Semgrep n'ont trouvé \textbf{aucune} des huit failles ; nos règles
  d'équipe en trouvent quatre, mais elles ont été écrites en connaissant les failles : ce biais surestime
  l'analyse statique. Sur un vrai projet, il faut des règles écrites avant de connaître les bugs.
  \item V2 et V3 sont invisibles pour l'analyse statique \emph{et} pour ZAP : ce sont des failles de
  logique métier et de contrôle d'accès. Seuls des tests qui connaissent les règles de l'application les
  couvrent. Le DAST complète le SAST, il ne remplace ni l'un ni l'autre la revue et les tests.
  \item Les outils produisent aussi du bruit : deux avertissements Semgrep et une alerte ZAP High étaient
  des faux positifs. Bloquer sans tri pousserait l'équipe à désactiver les contrôles.
\end{itemize}

\textbf{Erreurs commises, et ce qu'elles ont appris.}
\begin{itemize}
  \item Un fichier \code{trivy.yaml} ajouté pour accélérer les scans locaux excluait \code{node\_modules}
  \emph{aussi} dans le scan d'image : l'image de l'API passait avec la faille critique de lodash. Repéré en
  vérifiant un résultat trop beau, corrigé. Leçon : un contrôle qui passe doit aussi être vérifié.
  \item Les premiers compteurs de détection étaient en mémoire : avec deux répliques, trois sondes
  réparties 2 + 1 ne déclenchaient rien. Ils sont maintenant partagés en base.
  \item Le test V3 choisissait \og la commande d'un autre \fg{} d'après l'e-mail saisi ; ZAP avait créé
  des commandes avec des e-mails fantaisistes, et le test a échoué à tort. Il compare maintenant les
  comptes propriétaires.
  \item Le verrouillage a bloqué l'administrateur pendant une démonstration : c'est son rôle, mais la
  procédure de déblocage a été ajoutée au plan de réponse.
\end{itemize}

\textbf{Limites restantes.}
\begin{itemize}
  \item L'historique a été purgé et réécrit (Gitleaks : 0), mais les anciennes valeurs ont été publiques
  un temps : des copies peuvent subsister (clones, caches). C'est pourquoi la révocation passait d'abord.
  \item Environnement de test : pas de TLS, Secret Kubernetes non chiffré dans etcd, cluster d'un nœud ;
  sur kind, le trafic arrive avec l'adresse du nœud, ce qui affaiblit les règles par adresse.
  \item Pas de centralisation des journaux (SIEM) ni de tableau de bord ; pas de détection au niveau des
  conteneurs (Falco) ; images non signées (cosign) ; feuille Google Fonts sans empreinte SRI.
\end{itemize}

\textbf{Si c'était à refaire.} Écrire le modèle de menace et les tests de régression \emph{avant} le
code ; intégrer les tests métier au pipeline dès le premier jour ; ajouter TLS, la signature des images,
le chiffrement des secrets au repos et un SIEM.

\textbf{Déclaration d'usage d'outils d'IA.} Un assistant d'IA (Claude, Anthropic, via Claude Code) a été
utilisé pour produire une partie du code, des configurations, des scripts, de la documentation et de ce
rapport. Chaque résultat a été exécuté et vérifié par l'équipe sur ses propres instances (tests,
scanners, captures du pipeline) ; l'équipe reste responsable des choix et du contenu.

% ================================================================
\appendix
\section{Reproduire}
{\footnotesize
\begin{verbatim}
bash scripts/new-secrets.sh && docker compose up -d --build   # http://localhost:8081
bash scripts/k8s/up.sh && bash scripts/k8s/verify.sh           # cluster durci, 22 controles
BASE_URL=http://127.0.0.1:8090 npm --prefix backend run test:security
bash scripts/zap-scan.sh && bash scripts/demo/attack.sh       # ZAP, puis detection
\end{verbatim}}
Documents détaillés, dans \code{docs/} : \emph{vulnerabilities}, \emph{corrections}, \emph{pipeline-policy}, \emph{secrets}, \emph{kubernetes}, \emph{dast}, \emph{detection} et \emph{incident-response} (fichiers \code{.md}).

\end{document}
